Есть проблема поиска баланса между недообучением и переобучением.

Компромисс между смещением и дисперсией -- это свойство набора моделей предсказания, когда модели с меньшим
отклонением (low bias) от имеющихся данных имеют более высокую дисперсию (high variance) на новых данных (то есть
подвержены переобучению), и наоборот модели с меньшей дисперсией (low variance) от имеющихся данных имеют более
высокое смещение (high bias) на новых данных (это приводит к недообучению).
Притом всегда остаётся неустранимая случайность.
Ниже более подробное объяснение этого процесса.
Возникает конфликт при попытке одновременно минимизировать эти два источника ошибок, которые мешают алгоритмам
обучения с учителем делать обобщение за пределами тренировочного набора.

\subsection{Случайность в ошибке}
Случайность возникает
\begin{itemize}
    \item Из данных
    \item Из-за точности
    \item В процессе разбиения на тренировочный и тестовый наборы
    \item Вследствие человеческого фактора
    \item Пересборки тренировочного набора
\end{itemize}

\subsubsection{Интуиция каждой компоненты}

Итого,
\begin{enumerate}
    \item Random sampling in training procedure
    \subitem for Linear model -- initial weights
    \subitem for Tress -- feature and threshold sampling
    \subitem for SGD -- sampling of batches
    \item Computational errors due to limited precision arithmetics
    \item Target noise
    \subitem $\varepsilon$ in regression model
    \item Objects sampling -- crucial source of variance
    \subitem Train set sampling
    \subitem Validation set sampling
\end{enumerate}

Смещение -- это разница между средним прогнозом нашей модели и правильным значением, которое мы пытаемся предсказать.
Модель с высоким смещением уделяет очень мало внимания данным обучения и чрезмерно упрощает модель.
Это всегда приводит к высокой погрешности в обучающих и тестовых данных (недообучение).

Дисперсия -- это изменчивость прогноза модели для данной точки данных или значения, которое говорит нам о разбросе
наших данных.
Модель с высокой дисперсией уделяет большое внимание обучающим данным и не обобщается на данные, которые она раньше
не видела.
В результате такие модели очень хорошо работают с обучающими данными, но имеют высокую частоту ошибок в тестовых
данных (переобучение)

\subsection{Доказательство для MSE}

Рассмотрим задачу регрессии с квадратичной функцией потерь.
Представим также для простоты, что целевая переменная $y$ --- одномерная и выражается через переменную $x$ как:

\[ y = f(x) + \varepsilon, \]

где $f$ --- некоторая детерминированная функция, а $\varepsilon$ --- случайный шум со следующими свойствами:

\[ \mathbb{E}\,\varepsilon = 0, \quad \operatorname{Var}\varepsilon = \mathbb{E}\,\varepsilon^{2} = \sigma^{2} \]

В зависимости от природы данных, которые описывает эта зависимость, её представление в виде точной $f(x)$ и
случайной $\varepsilon$ может быть продиктовано по-разному

\begin{itemize}
    \item Данные на самом деле имеют случайный характер;
    \item измерительный прибор не может зафиксировать целевую переменную абсолютно точно;
    \item имеющихся признаков недостаточно, чтобы исчерпывающим образом описать объект, пользователя или событие.
\end{itemize}

Функция потерь на одном объекте $x$ равна
\[ \operatorname{MSE} = \bigl(y(x) - a(x)\bigr)^{2} \]

Однако знание значения MSE только на одном объекте не может дать нам общего понимания того, насколько хорошо
работает наш алгоритм.
Какие факторы мы бы хотели учесть при оценке качества алгоритма?
Например, то, что выход алгоритма на объекте $x$ зависит не только от самого этого объекта, но и от выборки $X$, на
которой алгоритм обучался:
\[ X = \bigl((x_{1},y_{1}), \dots, (x_{\ell}, y_{\ell})\bigr), \qquad a(x) = a(x, X) \]

Кроме того, значение $y$ на объекте $x$ зависит не только от $x$, но и от реализации шума в этой точке:
\[ y(x) = y(x, \varepsilon) \]

Наконец, измерять качество мы бы хотели на тестовых объектах $x$ --- тех, которые не встречались в обучающей
выборке, а тестовых объектов у нас в большинстве случаев более одного.
При включении всех вышеперечисленных источников случайности в рассмотрение логичной оценкой качества алгоритма $a$
кажется следующая величина:
\[ Q(a) = \mathbb{E}_{x}\, \mathbb{E}_{X,\varepsilon}\bigl[y(x,\varepsilon) - a(x,X)\bigr]^{2} \]

Внутреннее матожидание позволяет оценить качество работы алгоритма в одной тестовой точке $x$ в зависимости от
всевозможных реализаций $X$ и $\varepsilon$, а внешнее матожидание усредняет это качество по всем тестовым точкам.

\smallskip

\textbf{Замечание.}
Запись $\mathbb{E}_{X,\varepsilon}$ в общем случае обозначает взятие матожидания по совместному распределению $X$ и
$\varepsilon$.
Однако поскольку $X$ и $\varepsilon$ независимы, она равносильна последовательному взятию матожиданий по каждой из
переменных: $\mathbb{E}_{X,\varepsilon} = \mathbb{E}_{X}\,\mathbb{E}_{\varepsilon}$, но последний вариант выглядит
несколько более громоздко.

\medskip

Попробуем представить выражение для $Q(a)$ в более удобном для анализа виде.
Начнём с внутреннего матожидания:

\begin{align*}
    \mathbb{E}_{X,\varepsilon}\bigl[y(x,\varepsilon) - a(x,X)\bigr]^{2}
    &= \mathbb{E}_{X,\varepsilon}\bigl[f(x) + \varepsilon - a(x,X)\bigr]^{2} \\
    &= \mathbb{E}_{X,\varepsilon}\Bigl[
        \underbrace{(f(x) - a(x,X))^{2}}_{\text{не зависит от }\varepsilon}
        + 2\varepsilon\cdot\underbrace{(f(x) - a(x,X))}_{\text{множители независимы}}
        + \varepsilon^{2}
        \Bigr] \\
    &= \mathbb{E}_{X}\bigl[(f(x) - a(x,X))^{2}\bigr]
    + 2\underbrace{\mathbb{E}_{\varepsilon}[\varepsilon]}_{=0}\cdot\mathbb{E}_{X}\bigl(f(x) - a(x,X)\bigr)
    + \mathbb{E}_{\varepsilon}\,\varepsilon^{2} \\
    &= \mathbb{E}_{X}\bigl[(f(x) - a(x,X))^{2}\bigr] + \sigma^{2}.
\end{align*}

Из общего выражения для $Q(a)$ выделилась шумовая компонента $\sigma^{2}$.
Продолжим преобразования:

\begin{align*}
    \mathbb{E}_{X}\bigl[(f(x) - a(x,X))^{2}\bigr]
    &= \mathbb{E}_{X}\bigl[(f(x) - \mathbb{E}_{X}[a(x,X)] + \mathbb{E}_{X}[a(x,X)] - a(x,X))^{2}\bigr] \\
    &= \underbrace{\bigl(f(x) - \mathbb{E}_{X}[a(x,X)]\bigr)^{2}}_{\text{bias}_{X}^{2}\,a(x,X)}
    + \underbrace{\mathbb{E}_{X}\bigl[(a(x,X) - \mathbb{E}_{X}[a(x,X)])^{2}\bigr]}_{\operatorname{Var}_{X}[a(x,X)]} \\
    &\quad + 2\bigl(f(x) - \mathbb{E}_{X}[a(x,X)]\bigr)\cdot
    \underbrace{\mathbb{E}_{X}\bigl[\mathbb{E}_{X}[a(x,X)] - a(x,X)\bigr]}_{=0} \\
    &= \operatorname{bias}_{X}^{2}a(x,X) + \operatorname{Var}_{X}[a(x,X)].
\end{align*}

Таким образом, итоговое выражение для $Q(a)$ примет вид

\[ Q(a) = \mathbb{E}_{x}\,\mathbb{E}_{X,\varepsilon}\bigl[y(x,\varepsilon) - a(x,X)\bigr]^{2}
= \mathbb{E}_{x}\operatorname{bias}_{X}^{2}a(x,X)
+ \mathbb{E}_{x}\operatorname{Var}_{X}[a(x,X)]
+ \sigma^{2}, \]

Здесь
\[ \operatorname{bias}_{X}a(x,X) = f(x) - \mathbb{E}_{X}[a(x,X)] \] --- смещение предсказания алгоритма в точке $x$,
усреднённого по всем возможным обучающим выборкам, относительно истинной зависимости $f$;

\[ \operatorname{Var}_{X}[a(x,X)] = \mathbb{E}_{X}\bigl[a(x,X) - \mathbb{E}_{X}[a(x,X)]\bigr]^{2} \]
--- дисперсия (разброс) предсказаний алгоритма в зависимости от обучающей выборки $X$;

\[ \sigma^{2} = \mathbb{E}_{x}\,\mathbb{E}_{\varepsilon}\bigl[y(x,\varepsilon) - f(x)\bigr]^{2} \]
--- неустранимый шум в данных.

Смещение показывает, насколько хорошо с помощью данного алгоритма можно приблизить истинную зависимость $f$, а
разброс характеризует чувствительность алгоритма к изменениям в обучающей выборке.
Например, деревья маленькой глубины будут в большинстве случаев иметь высокое смещение и низкий разброс предсказаний,
так как они не могут слишком хорошо запомнить обучающую выборку.
А глубокие деревья, наоборот, могут безошибочно выучить обучающую выборку и потому будут иметь высокий разброс в
зависимости от выборки, однако их предсказания в среднем будут точнее.
На рисунке ниже приведены возможные случаи сочетания смещения и разброса для разных моделей

Синяя точка соответствует модели, обученной на некоторой обучающей выборке, а всего синих точек столько, сколько
было обучающих выборок.
Красный круг в центре области представляет ближайшую окрестность целевого значения.
Большое смещение соответствует тому, что модели в среднем не попадают в цель, а при большом разбросе модели могут
как делать точные предсказания, так и довольно сильно ошибаться.

\subsection{Компромисс bias–variance}

При построении моделей стремятся получить и низкий bias и адекватный variance

\subsection{Анализ основных моделей}

\begin{itemize}
    \item  kNN: low variance, average bias
    \item Linear models: very low variance, very high bias
    \item Trees: high variance, low bias for big trees.
    Depends on hyperparameters
\end{itemize}

\subsection{За пределами MSE}

Полученное нами разложение ошибки на три компоненты верно только для квадратичной функции потерь.
Для других функций потерь существуют более общие формы этого разложения с похожими по смыслу компонентами.
Это позволяет предполагать, что для большинства основных функций потерь имеется некоторое представление в виде
смещения, разброса и шума (хоть и, возможно, не в столь простой аддитивной форме)

\subsection{Double descent}

Double Descent -- феномен в машинном обучении, при котором ошибка теста модели сначала уменьшается, затем
увеличивается, а затем снова уменьшается по мере роста сложности модели.
Это контрастирует с традиционной U-образной кривой компромисса между предвзятостью и дисперсией, которая
предполагает, что ошибка теста должна монотонно увеличиваться с ростом сложности модели после определённого момента.

Явление Double Descent наблюдается в сложных моделях машинного обучения, например в глубоких нейронных сетях.

\underline{Механизм}

Процесс Double Descent включает три фазы
\begin{enumerate}
    \item Фаза недопереобучения -- по мере увеличения сложности модели ошибка теста уменьшается.
    Это характерно для моделей, которые не достаточно сложны, чтобы захватывать основные закономерности в данных.
    \item Фаза переобучения -- дальнейшее увеличение сложности модели приводит к увеличению ошибки теста.
    На этой фазе модели слишком сложны по сравнению с объёмом обучающих данных, захватывают шум как сигнал, что
    приводит к высокой дисперсии.
    \item Второй спуск -- по мере дальнейшего роста сложности модели (в режиме переобучения) ошибка теста снова
    начинает уменьшаться.
\end{enumerate}

\underline{Причины}

Явление Double Descent возникает из-за взаимодействия сложности модели, размера данных и шума.
По мере увеличения сложности модели она сначала переобучается шуму в данных, что увеличивает ошибку теста.
Однако по мере того, как модель становится ещё более сложной, она начинает интерполировать данные, меньше подгоняя
шум и больше -- сигнал, что снова уменьшает ошибку теста.